\begin{derivation}[从\eqnrefs{eq:compton-M}到\eqnrefs{eq:klein-nishina-SI}]
\begin{equation*}
 a\equiv p\cdot k,\qquad b\equiv p\cdot k',
\end{equation*}
并写
\begin{equation*}
 \mathcal K_{\rm BW}^{\mu\nu}
 =\gamma^\mu\frac{\slashed p+\slashed k+m_{\mathrm e} c}{2a}\gamma^\nu
 -\gamma^\nu\frac{\slashed p-\slashed k'+m_{\mathrm e} c}{2b}\gamma^\mu .
\end{equation*}
初态电子自旋和光子偏振各平均一次，末态求和。由于完整振幅已经满足 Ward 条件，偏振和中的规范项不贡献，因此
\begin{equation*}
 \overline{|\mathcal M_C|^2}
 =\frac{g_{\rm em}^4}{4}
 \Tr\!\left[(\slashed p'+m_{\mathrm e} c)\mathcal K_{\rm BW}^{\mu\nu}
 (\slashed p+m_{\mathrm e} c)\overline{\mathcal K}_{{\rm BW},\mu\nu}\right].
\end{equation*}
质量壳与动量守恒给出
\begin{equation*}
 k\cdot k'=a-b,\qquad
 p\cdot p'=m_{\mathrm e}^2c^2+a-b,
\end{equation*}
以及内部电子动量
\begin{equation*}
 r_s=p+k,\quad r_u=p-k',\quad
 r_s^2=m_{\mathrm e}^2c^2+2a,\quad r_u^2=m_{\mathrm e}^2c^2-2b.
\end{equation*}
把迹分成 $ss$、$uu$ 与两图干涉三部分。对 $ss$ 项，先用
\begin{equation*}
 \gamma^\nu(\slashed p+m_{\mathrm e} c)\gamma_\nu=-2\slashed p+4m_{\mathrm e} c
\end{equation*}
和
\begin{equation*}
 \slashed r\slashed p\slashed r=2(r\cdot p)\slashed r-r^2\slashed p
\end{equation*}
化简内部夹心，得到
\begin{align*}
 T_{ss}
 &=\Tr\!\left[(\slashed p'+m_{\mathrm e} c)\gamma^\mu(\slashed r_s+m_{\mathrm e} c)
 \gamma^\nu(\slashed p+m_{\mathrm e} c)\gamma_\nu(\slashed r_s+m_{\mathrm e} c)\gamma_\mu\right]\\
 &=32\left(m_{\mathrm e}^4c^4+m_{\mathrm e}^2c^2a+ab\right).
\end{align*}
连同平均因子和 $(2a)^{-2}$ 的传播子分母，
\begin{equation*}
 \overline{|\mathcal M_s|^2}
 =2g_{\rm em}^4\left(\frac ba+\frac{m_{\mathrm e}^2c^2}{a}+\frac{m_{\mathrm e}^4c^4}{a^2}\right).
\end{equation*}
$uu$ 项同样给出
\begin{equation*}
 \overline{|\mathcal M_u|^2}
 =2g_{\rm em}^4\left(\frac ab-\frac{m_{\mathrm e}^2c^2}{b}+\frac{m_{\mathrm e}^4c^4}{b^2}\right).
\end{equation*}
交叉迹在两种次序下相同，
\begin{equation*}
 T_{su}=T_{us}=16m_{\mathrm e}^2c^2(2m_{\mathrm e}^2c^2+a-b),
\end{equation*}
而正文振幅中的相对负号使干涉贡献成为
\begin{equation*}
 \overline{2\Re(\mathcal M_s\mathcal M_u^*)}
 =2g_{\rm em}^4\left[
 m_{\mathrm e}^2c^2\left(\frac1a-\frac1b\right)
 -\frac{2m_{\mathrm e}^4c^4}{ab}\right].
\end{equation*}
三部分相加并整理同次幂的 $1/a,1/b$，得到正文\eqnrefs{eq:compton-M2-unpolarized-dp73}。

电子静止系中
\begin{equation*}
 p^\mu=(m_{\mathrm e} c,\mathbf0),\qquad
 a=m_{\mathrm e}\hbar\omega,\qquad b=m_{\mathrm e}\hbar\omega'.
\end{equation*}
利用正文反冲关系~\eqnrefs{eq:compton-energy-ratio}，\eqnrefs{eq:compton-M2-unpolarized-dp73} 化为
\begin{equation*}
 \overline{|\mathcal M_C|^2}
 =2g_{\rm em}^4\left(
 \frac\omega{\omega'}+\frac{\omega'}\omega-\sin^2\theta\right).
\end{equation*}

相空间从 Lorentz 不变二体测度开始。末态电子能量记为
\begin{equation*}
 E_e'=\sqrt{m_{\mathrm e}^2c^4+c^2|\mathbf p'|^2}.
\end{equation*}
用空间 $\delta$ 函数消去 $\mathbf p'$ 后，
\begin{equation*}
 \dd\mathfrak P_2
 =\frac{\hbar^2\omega'}{16\pi^2E_e'}
 \dd\omega'\dd\Omega\,
 \delta\!\left(m_{\mathrm e} c^2+\hbar\omega-E_e'-\hbar\omega'\right).
\end{equation*}
令括号内的能量守恒函数为 $f(\omega')$。由
\begin{equation*}
 \mathbf p'=\frac{\hbar}{c}
 (\omega\hat{\mathbf k}-\omega'\hat{\mathbf k}')
\end{equation*}
可得物理解处
\begin{align*}
 |f'(\omega')|
 &=\frac{\hbar}{E_e'}
 \left[m_{\mathrm e} c^2+\hbar\omega(1-\cos\theta)\right]\\
 &=\frac{\hbar m_{\mathrm e} c^2}{E_e'}\frac\omega{\omega'}.
\end{align*}
因此用 $\delta[f]=\delta(\omega'-\omega'_*)/|f'(\omega'_*)|$ 完成能量积分后，
\begin{equation*}
 \frac{\dd\mathfrak P_2}{\dd\Omega}
 =\frac{\hbar}{16\pi^2m_{\mathrm e} c^2}\frac{\omega'^2}{\omega}.
\end{equation*}
正文\eqnrefs{eq:compton-invariant-flux-r14} 给出入射通量
$\mathcal I_{e\gamma}=m_{\mathrm e}\hbar\omega$。代入 SI 约化振幅的截面母式
\begin{equation*}
 \dd\sigma
 =\frac{\hbar^2\overline{|\mathcal M|^2}}{4\mathcal I_{12}}
 \dd\mathfrak P_2
\end{equation*}
得到
\begin{equation*}
 \frac{\dd\sigma_C}{\dd\Omega}
 =\frac{\hbar^2}{64\pi^2m_{\mathrm e}^2c^2}
 \left(\frac{\omega'}\omega\right)^2
 \overline{|\mathcal M_C|^2}.
\end{equation*}
最后使用 $g_{\rm em}^2=4\pi\alpha$ 与
$r_e=\alpha\hbar/(m_{\mathrm e} c)$，便得到正文\eqnrefs{eq:klein-nishina-SI}。
\end{derivation}

\begin{derivation}[从\eqnrefs{eq:compton-M}到\eqnrefs{eq:ee-gammagamma-M-r68}]
Dirac 场展开含

\begin{equation*}
 \hat\psi(x)\supset b_s(\mathbf p)u_s(p)\ee^{-\ii p\cdot x/\hbar}
 +d_s^\dagger(\mathbf p)v_s(p)\ee^{+\ii p\cdot x/\hbar}.
\end{equation*}
入射电子外腿从正频项取得 $u(p)\ee^{-\ii p\cdot x/\hbar}$；同一条费米子线改解释为出射正电子时，收缩从负频项取得 $v(p)\ee^{+\ii p\cdot x/\hbar}$。顶角的 Fourier 积分因此对应四动量解析延拓 $p\mapsto-p$，同时把 $u$ 外线替换为相应 $v$ 外线。交叉对称性正是同一时序 Green 函数在不同质量壳边界值之间的解析延拓。

对
\begin{equation*}
 e^-(p_1)+e^+(p_2)\to\gamma(k_1)+\gamma(k_2),
\end{equation*}
两个光子顶角的两种次序给出
\begin{equation*}
 {
 \mathcal M_{\gamma\gamma}
 =-g_{\rm em}^2\bar v(p_2)\left[
 \slashed\varepsilon_2^*
 \frac{\slashed p_1-\slashed k_1+m_{\mathrm e} c}{(p_1-k_1)^2-m_{\mathrm e}^2c^2}
 \slashed\varepsilon_1^*
 +\slashed\varepsilon_1^*
 \frac{\slashed p_1-\slashed k_2+m_{\mathrm e} c}{(p_1-k_2)^2-m_{\mathrm e}^2c^2}
 \slashed\varepsilon_2^*
 \right]u(p_1).}
 \end{equation*}
\end{derivation}


\begin{derivation}[从\eqnrefs{eq:BW-amplitude-r68}到\eqnrefs{eq:BW-dsigma-si-dp49}]
正文\eqnrefs{eq:BW-amplitude-r68} 给出两张电子交换图的相干和。先完成自旋--偏振求和和二体相空间，只推到正文角分布；总截面的立体角积分留给下一条推导。

定义
\begin{equation*}
 s=(k_1+k_2)^2,\qquad
 t=(k_1-p_1)^2,\qquad
 u=(k_2-p_1)^2,
\end{equation*}
并使用正文\eqnrefs{eq:BW-stu-si-r12} 的 $s+t+u=2m_{\mathrm e}^2c^2$。质心系末态速率参数为
\begin{equation*}
 \beta_{\rm pair}=\sqrt{1-\frac{4m_{\mathrm e}^2c^2}{s}},
\end{equation*}
因此
\begin{equation*}
 |\mathbf k_1|=\frac{\sqrt s}{2},\qquad
 |\mathbf p_1|=\frac{\sqrt s}{2}\beta_{\rm pair}.
\end{equation*}
若 $\theta$ 是 $\mathbf k_1$ 与 $\mathbf p_1$ 的夹角，则
\begin{align*}
 t-m_{\mathrm e}^2c^2&=-\frac{s}{2}(1-\beta_{\rm pair}\cos\theta),\\
 u-m_{\mathrm e}^2c^2&=-\frac{s}{2}(1+\beta_{\rm pair}\cos\theta).
\end{align*}

Breit--Wheeler 与 Compton 属于同一个截肢四点函数的交叉道，因此八伽马迹的 Lorentz 不变量结构已经包含在\eqnrefs{eq:compton-M2-unpolarized-dp73} 中。对该结果作交叉解析延拓后，只需重新确定交叉后的 Mandelstam 不变量以及费米子自旋和所带的额外负号。

把 Compton 过程
\begin{equation*}
 e^-(p)+\gamma(k)\to e^-(p')+\gamma(k')
\end{equation*}
交叉成
\begin{equation*}
 \gamma(k_1)+\gamma(k_2)\to e^-(p_1)+e^+(p_2)
\end{equation*}
时取
\begin{equation*}
 p\to-p_2,
 \qquad
 k\to k_1,
 \qquad
 p'\to p_1,
 \qquad
 k'\to-k_2.
\end{equation*}
令
\begin{equation*}
 T\equiv t-m_{\mathrm e}^2c^2,
 \qquad
 U\equiv u-m_{\mathrm e}^2c^2.
\end{equation*}
则 Compton 中的两个不变量
\begin{equation*}
 a=p\cdot k,
 \qquad
 b=p\cdot k'
\end{equation*}
分别变成
\begin{align*}
 a&\to-p_2\cdot k_1
 =\frac{u-m_{\mathrm e}^2c^2}{2}=\frac U2,\\
 b&\to p_2\cdot k_2
 =-\frac{t-m_{\mathrm e}^2c^2}{2}=-\frac T2.
\end{align*}
第二行使用了四动量守恒 $k_2-p_2=p_1-k_1$。

平方振幅的交叉还必须处理费米子完备关系。Compton 初态电子自旋和含
\begin{equation*}
 \sum_su_s(p)\bar u_s(p)=\slashed p+m_{\mathrm e} c.
\end{equation*}
解析延拓 $p\to-p_2$ 后，
\begin{equation*}
 \slashed p+m_{\mathrm e} c
 \longrightarrow
 -\slashed p_2+m_{\mathrm e} c
 =-\bigl(\slashed p_2-m_{\mathrm e} c\bigr),
\end{equation*}
而括号内正是末态正电子的完备关系
\begin{equation*}
 \sum_sv_s(p_2)\bar v_s(p_2)=\slashed p_2-m_{\mathrm e} c.
\end{equation*}
因此交叉后的未偏振迹相对于对 Compton 不变量函数的直接代入多出一个整体负号：
\begin{equation*}
 \overline{|\mathcal M_{\rm BW}|^2}
 =-\,\overline{|\mathcal M_C|^2}
 \bigg|_{a=U/2,\,b=-T/2}.
\end{equation*}
把\eqnrefs{eq:compton-M2-unpolarized-dp73} 逐项代入：
\begin{align*}
 \overline{|\mathcal M_{\rm BW}|^2}
 =2g_{\rm em}^4\Bigg[
 &\frac UT+\frac TU
 -4m_{\mathrm e}^2c^2\left(\frac1U+\frac1T\right)\\
 &-4m_{\mathrm e}^4c^4\left(\frac1U+\frac1T\right)^2
 \Bigg].
\end{align*}
由
\begin{equation*}
 s+t+u=2m_{\mathrm e}^2c^2
 \quad\Longrightarrow\quad
 T+U=-s,
\end{equation*}
可把它等价写成
\begin{equation*}
 \overline{|\mathcal M_{\rm BW}|^2}
 =2g_{\rm em}^4\left[
 \frac UT+\frac TU
 +\frac{4m_{\mathrm e}^2c^2s}{TU}
 -\frac{4m_{\mathrm e}^4c^4s^2}{T^2U^2}
 \right].
\end{equation*}
这已经是完全 Lorentz 不变的 Breit--Wheeler 树级平方振幅。

在质心系中
\begin{equation*}
 T=-\frac s2(1-\beta_{\rm pair}\cos\theta),
 \qquad
 U=-\frac s2(1+\beta_{\rm pair}\cos\theta),
\end{equation*}
并且
\begin{equation*}
 m_{\mathrm e}^2c^2=\frac s4(1-\beta_{\rm pair}^2).
\end{equation*}
因此
\begin{align*}
 TU
 &=\frac{s^2}{4}(1-\beta_{\rm pair}^2\cos^2\theta),\\
 T^2+U^2
 &=\frac{s^2}{2}(1+\beta_{\rm pair}^2\cos^2\theta).
\end{align*}
把这些量代入上一式并通分，分子逐项为
\begin{align*}
 &2(1+\beta_{\rm pair}^2\cos^2\theta)
 (1-\beta_{\rm pair}^2\cos^2\theta)\\
 &\quad+4(1-\beta_{\rm pair}^2)
 (1-\beta_{\rm pair}^2\cos^2\theta)
 -4(1-\beta_{\rm pair}^2)^2.
\end{align*}
展开并合并同次幂后得到
\begin{equation*}
 2\left[
 1-\beta_{\rm pair}^4\cos^4\theta
 +2\beta_{\rm pair}^2(1-\beta_{\rm pair}^2)\sin^2\theta
 \right].
\end{equation*}
所以
\begin{equation*}
 \boxed{
 \overline{|\mathcal M_{\rm BW}|^2}
 =4g_{\rm em}^4
 \frac{
 1-\beta_{\rm pair}^4\cos^4\theta
 +2\beta_{\rm pair}^2(1-\beta_{\rm pair}^2)\sin^2\theta
 }{(1-\beta_{\rm pair}^2\cos^2\theta)^2}.}
\end{equation*}
这样新出现的数学内容只有交叉映射及其符号，而重复的 Dirac 迹由已经完成的 Compton 推导承担。
二体截面母\eqnrefs{eq:2to2-cross-section-master-si-r9} 给出
\begin{equation*}
 \frac{\dd\sigma_{\rm BW}}{\dd\Omega}
 =\frac{\hbar^2\beta_{\rm pair}}{64\pi^2s}
 \overline{|\mathcal M_{\rm BW}|^2}.
\end{equation*}
最后代入
\begin{equation*}
 s=\frac{4m_{\mathrm e}^2c^2}{1-\beta_{\rm pair}^2},\qquad
 g_{\rm em}^4=(4\pi\alpha)^2,\qquad
 r_e=\frac{\alpha\hbar}{m_{\mathrm e} c},
\end{equation*}
便得到正文\eqnrefs{eq:BW-dsigma-si-dp49}。
\end{derivation}


\begin{derivation}[从\eqnrefs{eq:vacuum-polarization-loop}到\eqnrefs{eq:Pi-transverse-r12}]
记约化自由电子传播子为
\begin{equation*}
 S_0(\ell)=\frac{\slashed\ell+m_{\mathrm e} c}{\ell^2-m_{\mathrm e}^2c^2+\ii0^+},
 \qquad
 S_{0}^{-1}(\ell)=\slashed\ell-m_{\mathrm e} c .
\end{equation*}
由
\begin{equation*}
 \slashed q=S_0^{-1}(\ell+q)-S_0^{-1}(\ell)
\end{equation*}
收缩正文\eqnrefs{eq:vacuum-polarization-loop} 的外部动量，并把两个传播子的显式 $\ii$ 因子保留到底。利用迹的循环性得
\begin{align*}
 \ii q_\mu\Pi^{\mu\nu}(q)
 &=-g_{\rm em}^2
 \int_\ell^{(d)}
 \Tr\!\left[
 S_0(\ell+q)\slashed q\,S_0(\ell)\gamma^\nu
 \right]\\
 &=-g_{\rm em}^2
 \int_\ell^{(d)}
 \Tr\!\left[
 \bigl(S_0(\ell)-S_0(\ell+q)\bigr)\gamma^\nu
 \right].
\end{align*}
维数正则化保持圈动量平移不变；第二项作 \(\ell\mapsto\ell-q\) 后与第一项相消，因此
\begin{equation*}
 q_\mu\Pi^{\mu\nu}(q)=0.
\end{equation*}
真空 Lorentz 协变性只允许
\begin{equation*}
 \Pi^{\mu\nu}(q)=A(q^2)\eta^{\mu\nu}+B(q^2)q^\mu q^\nu .
\end{equation*}
与横向条件联立给 \(A=-q^2B\)，故得到正文\eqnrefs{eq:Pi-transverse-r12}。
\end{derivation}

\begin{derivation}[从\eqnrefs{eq:vacuum-polarization-loop}到\eqnrefs{eq:Pi-renormalized-r12}]
这里把一圈真空极化中最容易被 ``标准积分'' 掩盖的三步完整展开：先做 Dirac 迹，再用维数正则化中的分部积分恒等式把张量积分化成横向结构，最后才施加 $\Pi_R(0)=0$ 的电荷重整化条件。

记
\[
 A=\ell^2-m_{\mathrm e}^2c^2+\ii0^+,
 \qquad
 B=(\ell+q)^2-m_{\mathrm e}^2c^2+\ii0^+ .
\]
两分母参数化为
\[
 \frac1{AB}=\int_0^1\dd x\,
 \frac1{[xA+(1-x)B]^2}.
\]
令
\[
 L\equiv \ell+(1-x)q,
 \qquad
 \Delta_\Pi\equiv m_{\mathrm e}^2c^2-x(1-x)q^2,
\]
则
\[
 xA+(1-x)B=L^2-\Delta_\Pi+\ii0^+.
\]

闭合电子圈分子中的 Dirac 迹先按四伽马迹恒等式计算：
\begin{align*}
N^{\mu\nu}
&\equiv
\Tr\!\left[
\gamma^\mu(\slashed\ell+m_{\mathrm e} c)
\gamma^\nu(\slashed\ell+\slashed q+m_{\mathrm e} c)
\right]\\
&=4\left[
\ell^\mu(\ell+q)^\nu
+\ell^\nu(\ell+q)^\mu
-\eta^{\mu\nu}
\bigl(\ell\!\cdot(\ell+q)-m_{\mathrm e}^2c^2\bigr)
\right].
\end{align*}
代入
$\ell=L-(1-x)q$、$\ell+q=L+xq$。关于 $L$ 的奇次项在平移不变的维数正则化积分中为零，因此有效偶次分子为
\begin{align*}
N_{\rm even}^{\mu\nu}
=4\Bigl[
&2L^\mu L^\nu-\eta^{\mu\nu}L^2
+\eta^{\mu\nu}m_{\mathrm e}^2c^2\\
&+x(1-x)
\bigl(q^2\eta^{\mu\nu}-2q^\mu q^\nu\bigr)
\Bigr].
\end{align*}

此时不能只说 ``由对称性令 $L^\mu L^\nu\to\eta^{\mu\nu}L^2/d$'' 后跳到答案；真正保证横向性的关键是同一调节下的积分分部恒等式
\begin{align*}
0
&=\int_L^{(d)}
\frac{\partial}{\partial L_\mu}
\left[
\frac{L^\nu}{L^2-\Delta_\Pi+\ii0^+}
\right]\\
&=\eta^{\mu\nu}
\int_L^{(d)}\frac1{L^2-\Delta_\Pi+\ii0^+}
-2\int_L^{(d)}
\frac{L^\mu L^\nu}
{(L^2-\Delta_\Pi+\ii0^+)^2}.
\end{align*}
所以
\[
2\int_L^{(d)}
\frac{L^\mu L^\nu}{(L^2-\Delta_\Pi+\ii0^+)^2}
=\eta^{\mu\nu}
\int_L^{(d)}\frac1{L^2-\Delta_\Pi+\ii0^+}.
\]
另一方面
\[
\frac{L^2}{(L^2-\Delta_\Pi)^2}
=\frac1{L^2-\Delta_\Pi}
+\frac{\Delta_\Pi}{(L^2-\Delta_\Pi)^2}.
\]
因此前三项的积分恰好化成
\begin{align*}
&\int_L^{(d)}
\frac{2L^\mu L^\nu-\eta^{\mu\nu}L^2
+\eta^{\mu\nu}m_{\mathrm e}^2c^2}
{(L^2-\Delta_\Pi+\ii0^+)^2}\\
&\qquad
=\eta^{\mu\nu}
(m_{\mathrm e}^2c^2-\Delta_\Pi)
I_2(\Delta_\Pi)\\
&\qquad
=x(1-x)q^2\eta^{\mu\nu}I_2(\Delta_\Pi),
\end{align*}
其中
\[
I_2(\Delta)
\equiv
\int_L^{(d)}\frac1{(L^2-\Delta+\ii0^+)^2}.
\]
再与显式的 $q^\mu q^\nu$ 项合并，整个分子积分成为
\begin{align*}
\int_L^{(d)}
\frac{N_{\rm even}^{\mu\nu}}
{(L^2-\Delta_\Pi+\ii0^+)^2}
&=-8x(1-x)
(q^\mu q^\nu-q^2\eta^{\mu\nu})
I_2(\Delta_\Pi).
\end{align*}
横向张量因此不是事后猜出的结构，而是在同一个维数正则化积分中由分部积分恒等式直接产生。

把正文\eqnrefs{eq:vacuum-polarization-loop} 的整体因子恢复，并使用 $g_{\rm em}^2=4\pi\alpha$ 与正文\eqnrefs{eq:I2-expanded-r12}，即可从
\[
\Pi^{\mu\nu}(q)
=(q^\mu q^\nu-q^2\eta^{\mu\nu})\Pi(q^2)
\]
读出
\begin{align*}
\Pi(q^2)
=\frac{2\alpha}{\pi}
\int_0^1\dd x\,x(1-x)
\left[
\frac1\epsilon-\gamma_E+\ln4\pi
+\ln\frac{\mu_{\rm DR}^2}
{\Delta_\Pi(x,q^2)-\ii0^+}
\right]
+\order(\epsilon).
\end{align*}
在 $q^2=0$ 时，$\Delta_\Pi(x,0)=m_{\mathrm e}^2c^2$，所以同一表达式给
\begin{align*}
\Pi(0)
=\frac{2\alpha}{\pi}
\int_0^1\dd x\,x(1-x)
\left[
\frac1\epsilon-\gamma_E+\ln4\pi
+\ln\frac{\mu_{\rm DR}^2}{m_{\mathrm e}^2c^2}
\right].
\end{align*}
两式相减时，紫外极点、$\overline{\rm MS}$ 方案常数和任意重整化尺度同时消失：
\begin{align*}
\Pi_R(q^2)
&\equiv\Pi(q^2)-\Pi(0)\\
&=\frac{2\alpha}{\pi}
\int_0^1\dd x\,x(1-x)
\ln\frac{m_{\mathrm e}^2c^2}
{m_{\mathrm e}^2c^2-q^2x(1-x)-\ii0^+},
\end{align*}
这正是正文\eqnrefs{eq:Pi-renormalized-r12}。这里 $\Pi_R(0)=0$ 不是额外代数技巧，而是在壳电荷定义：零动量传递处测得的电荷被定义为 $e$，所有真空极化修正从有限的动量依赖部分开始。
\end{derivation}


\begin{derivation}[从\eqnrefs{eq:qed-dyson-transverse-r68}到\eqnrefs{eq:qed-alpha-effective-dp7}]
守恒外部电流只耦合光子传播子的横向部分。提出共同的 Lorentz 投影算符与裸分母后，一次已重整化真空极化插入给传播核乘上无量纲因子。Dyson 方程写成

\begin{equation*}
 D_T(q)=D_{0,T}(q)+D_{0,T}(q)\,\Pi_T(q)\,D_T(q).
 \end{equation*}
把第二项移到左边并提取 $D_T$：
\begin{align*}
 \left[D_{0,T}^{-1}(q)-\Pi_T(q)\right]D_T(q)&=1,\\
 D_T(q)&=\left[D_{0,T}^{-1}(q)-\Pi_T(q)\right]^{-1}.
\end{align*}
用前面定义的无量纲已重整化标量 $\Pi_R$ 参数化以后，空间样 $q^2=-Q^2$ 的横向传播子相对于裸传播子多出
\begin{equation*}
 {
 D_T(-Q^2)=\frac{D_{0,T}(-Q^2)}{1+\Pi_R(-Q^2)}.}
 \end{equation*}
因为树级电流--电流振幅只含 ``两个 QED 顶角乘一条光子传播子''，可以把同一个修正等效地吸收到耦合中，定义一圈-improved 有效 fine-结构常数
\begin{equation*}
 {
 \alpha_{\rm eff}(Q)
 \equiv
 \frac{\alpha}{1+\Pi_R(-Q^2)}.}
 \end{equation*}
\eqnrefs{eq:qed-alpha-effective-dp7} 只重求和反复出现的一圈单粒子不可约真空极化核；更高圈单粒子不可约图会继续修正 $\beta$ 函数，因此该表达式是一圈精度下的运行耦合。
\end{derivation}

\begin{derivation}[从\eqnrefs{eq:Pi-Delta-r12}到\eqnrefs{eq:ImPi-r12}]
当时间样 $q^2=s>4m_{\mathrm e}^2c^2$ 时，\eqnrefs{eq:Pi-Delta-r12} 在一段 Feynman-参数区间内变负。方程

\begin{equation*}
 m_{\mathrm e}^2c^2-sx(1-x)=0
\end{equation*}
等价于
\begin{equation*}
 x^2-x+\frac{m_{\mathrm e}^2c^2}{s}=0,
\end{equation*}
根为
\begin{equation*}
 x_\pm=\frac12(1\pm\beta_{\rm pair}),
 \qquad
 \beta_{\rm pair}\equiv\sqrt{1-\frac{4m_{\mathrm e}^2c^2}{s}}.
\end{equation*}
对 $x_-<x<x_+$，分母的实部为负。由主支对数
\begin{equation*}
 \ln(-A-\ii0^+)=\ln A-\ii\pi,
 \qquad A>0,
\end{equation*}
而 \eqnrefs{eq:Pi-renormalized-r12} 中的对数是分子除以该负数，所以
\begin{equation*}
 \ln\frac{m_{\mathrm e}^2c^2}{-A-\ii0^+}
 =\ln\frac{m_{\mathrm e}^2c^2}{A}+\ii\pi.
\end{equation*}
于是
\begin{align*}
 \operatorname{Im}\Pi_R(s)
 &=2\alpha\int_{x_-}^{x_+}\dd x\,x(1-x),\\
 &=2\alpha\left[
 \frac{x^2}{2}-\frac{x^3}{3}
 \right]_{x_-}^{x_+}.
\end{align*}
利用 $x_++x_-=1$、$x_+-x_-=\beta_{\rm pair}$ 与 $x_+x_-=m_{\mathrm e}^2c^2/s$，可化为
\begin{equation*}
 {
 \operatorname{Im}\Pi_R(s)
 =\frac{\alpha}{3}\,
 \beta_{\rm pair}\left(1+\frac{2m_{\mathrm e}^2c^2}{s}\right).}
 \end{equation*}
\end{derivation}

\begin{derivation}[从\eqnrefs{eq:selfenergy-start-r12}到\eqnrefs{eq:selfenergy-uv-r68}]
真空中除 $p^\mu$ 外没有其他优先四矢量。QED 保持宇称，自能是自旋空间中的 $4\times4$ 矩阵；Lorentz 协变性与宇称因此排除独立的 $\gamma^5$、$\gamma^\mu\gamma^5$ 与 $\sigma^{\mu\nu}$ 项，一般形式只能是

\begin{equation*}
 {
 \Sigma(\slashed p)
 =A(p^2)\slashed p+B(p^2)m_{\mathrm e} c.}
 \end{equation*}
Lorentz 协变性与 Ward 恒等式先限制顶角的允许结构，一圈积分再显式给出相应形状因子。Feynman 规范中，约化图规则给
\begin{equation*}
 \Sigma(\slashed p)
 =-\ii g_{\rm em}^2
 \int_{\ell}^{(d)}
 \frac{\gamma^\mu(\slashed p-\slashed\ell+m_{\mathrm e} c)\gamma_\mu}
 {[\ell^2+\ii0^+][(p-\ell)^2-m_{\mathrm e}^2c^2+\ii0^+]}.
 \end{equation*}
用相应的 $d$ 维 Clifford 恒等式
\begin{equation*}
 \gamma^\mu\gamma^\alpha\gamma_\mu=(2-d)\gamma^\alpha,
 \qquad
 \gamma^\mu\gamma_\mu=d\id(),
\end{equation*}
得到
\begin{equation*}
 \gamma^\mu(\slashed p-\slashed\ell+m_{\mathrm e} c)\gamma_\mu
 =(2-d)(\slashed p-\slashed\ell)+dm_{\mathrm e} c.
 \end{equation*}
用 \eqnrefs{eq:feynman-parameter-proof-r9} 合并两个分母：
\begin{align*}
 &x[(p-\ell)^2-m_{\mathrm e}^2c^2]+(1-x)\ell^2\\
 &\qquad=(\ell-xp)^2-\Delta_\Sigma(x,p^2),
\end{align*}
其中
\begin{equation*}
 {
 \Delta_\Sigma(x,p^2)
 =xm_{\mathrm e}^2c^2-x(1-x)p^2.}
 \end{equation*}
令 $L=\ell-xp$，则
\begin{equation*}
 \slashed p-\slashed\ell=(1-x)\slashed p-\slashed L.
\end{equation*}
关于 $L$ 为奇函数的 $\slashed L$ 项在对称积分中为零，因此
\begin{align*}
 \Sigma(\slashed p)
 =-\ii g_{\rm em}^2
 \int_0^1\dd x\,
 \bigl[(2-d)(1-x)\slashed p+dm_{\mathrm e} c\bigr]
 I_2(\Delta_\Sigma).
 \end{align*}
取 $d=4-2\epsilon$，只保留 \eqnrefs{eq:I2-expanded-r12} 的紫外极点：
\begin{align*}
 \Sigma_{\rm UV}(\slashed p)
 &=\frac{g_{\rm em}^2}{(4\pi)^2}\frac1\epsilon
 \int_0^1\dd x
 \left[-2(1-x)\slashed p+4m_{\mathrm e} c\right],\\
 &=\frac{\alpha}{4\pi}\frac1\epsilon
 \left(-\slashed p+4m_{\mathrm e} c\right).
 \end{align*}
因此电子场强反项与质量反项必须分别吸收 $\slashed p$ 和 $m_{\mathrm e} c$ 两种独立结构；Lorentz 协变的二点函数中没有第三个独立反项结构可供调节。
\end{derivation}

\begin{derivation}[从\eqnrefs{eq:once-subtracted-single-cut-dp9}到\eqnrefs{eq:Pi-spacelike-r12}]
已经由 Feynman-参数积分得到

\begin{equation*}
 \operatorname{Im}\Pi_R(s)
 =\frac{\alpha}{3}
 \beta_{\rm pair}(s)\left(1+\frac{2m_{\mathrm e}^2c^2}{s}\right),
 \qquad
 \beta_{\rm pair}(s)\equiv\sqrt{1-\frac{4m_{\mathrm e}^2c^2}{s}},
 \end{equation*}
其支撑从 $s_0=4m_{\mathrm e}^2c^2$ 开始。该阈值由光子单粒子不可约二点函数的 Cutkosky 切割决定：费米子与反费米子两条内部线同时满足
\begin{equation*}
 p_-^2=p_+^2=m_{\mathrm e}^2c^2,
 \qquad
 p_-^0>0,\quad p_+^0>0,
 \qquad
 q=p_-+p_+,
\end{equation*}
因此只有虚光子的 Lorentz 不变质量达到 $2m_{\mathrm e} c$ 后，真实 $e^-e^+$ 在壳相空间才打开。

在壳电荷重整化已固定 $\Pi_R(0)=0$。因此直接使用通用一次减法公式 \eqnrefs{eq:once-subtracted-single-cut-dp9}：
\begin{equation*}
 {
 \Pi_R(q^2)
 =\frac{q^2}{\pi}
 \int_{4m_{\mathrm e}^2c^2}^{\infty}\dd s\,
 \frac{\operatorname{Im}\Pi_R(s)}
 {s\,[s-q^2-\ii0^+]}.}
 \end{equation*}
这里减法常数已由物理重整化条件固定；割线上的谱权重则由幺正性决定，两者承担不同信息。

对空间样 $q^2=-Q^2$，$Q^2>0$，
\begin{equation*}
 \Pi_R(-Q^2)
 =-\frac{Q^2}{\pi}
 \int_{4m_{\mathrm e}^2c^2}^{\infty}\dd s\,
 \frac{\operatorname{Im}\Pi_R(s)}{s(s+Q^2)}<0,
 \end{equation*}
这与直接由 Feynman 参数得到的 \eqnrefs{eq:Pi-spacelike-r12} 符号一致。
\end{derivation}

\begin{derivation}[从\eqnrefs{eq:qed-vacpol-spacelike-dispersion-dp9}到\eqnrefs{eq:qed-pi-smallQ-dp7}]
令 $a\equiv m_{\mathrm e}^2c^2$ 只在本段积分中作紧凑记号。由 \eqnrefs{eq:qed-vacpol-spacelike-dispersion-dp9}，当 $Q^2\ll a$，

\begin{align*}
 \Pi_R(-Q^2)
 &=-\frac{Q^2}{\pi}
 \int_{4a}^{\infty}\dd s\,
 \frac{\operatorname{Im}\Pi_R(s)}{s^2}
 +\order\!\left(\frac{Q^4}{a^2}\right).
 \end{align*}
把 \eqnrefs{eq:qed-vacpol-spectral-weight-dp9} 代入，并作变量
\begin{equation*}
 y\equiv\sqrt{1-\frac{4a}{s}},
 \qquad
 s=\frac{4a}{1-y^2},
 \qquad
 \dd s=\frac{8ay}{(1-y^2)^2}\dd y,
\end{equation*}
其中 $s:4a\to\infty$ 对应 $y:0\to1$。于是
\begin{align*}
 \int_{4a}^{\infty}\dd s\,
 \frac{\operatorname{Im}\Pi_R(s)}{s^2}
 &=\frac{\alpha}{3}
 \int_0^1\dd y\,
 \frac{y^2(3-y^2)}{4a},\\
 &=\frac{\alpha}{12a}
 \left[y^3-\frac{y^5}{5}\right]_0^1,\\
 &=\frac{\alpha}{15a}.
\end{align*}
因此
\begin{equation*}
 {
 \Pi_R(-Q^2)
 =-\frac{\alpha}{15\pi}
 \frac{Q^2}{m_{\mathrm e}^2c^2}
 +\order\!\left(\frac{Q^4}{m_{\mathrm e}^4c^4}\right),}
 \end{equation*}
精确恢复直接圈积分的 \eqnrefs{eq:qed-pi-smallQ-dp7}。这是一项非平凡交叉检查：一个系数由欧几里得/空间样 Feynman 参数展开得到，另一个只由时间样在壳对-产生谱权重得到。
\end{derivation}

\begin{derivation}[从\eqnrefs{eq:selfenergy-start-r12}到\eqnrefs{eq:selfenergy-onshell-no-decay-r68}]
一圈电子自能 \eqnrefs{eq:selfenergy-start-r12} 的两个内部分母对应一条光子线与一条电子线。若先以 $m_\gamma>0$ 作为光子质量的红外调节参数，两体切割的运动学阈值为

\begin{equation*}
 p^2\ge (m_{\mathrm e}+m_\gamma)^2c^2.
 \end{equation*}
因此调节参数存在时，物理电子极点 $p^2=m_{\mathrm e}^2c^2$ 与连续谱阈值分离，精确传播子具有通常的孤立极点；对当前场变量，其 LSZ 极点权重记为 $Z_{\rm pole}^{(e)}$，可由极点邻域的逆传播子导数求出。若随后采用在壳重整化，则场强反项中的 $Z_2$ 被选择来把重整化传播子的这一物理极点权重归一化为 $1$。

现在令 $m_\gamma\to0$。阈值向 $p^2=m_{\mathrm e}^2c^2$ 靠近，但这不意味着自由电子可以发出一个有限能量实光子。若假设
\begin{equation*}
 p\to p'+k,
 \qquad
 p^2=p'^2=m_{\mathrm e}^2c^2,
 \qquad
 k^2=0,
\end{equation*}
四动量守恒 $p=p'+k$ 两边平方：
\begin{align*}
 p^2
 &=(p'+k)^2,\\
 m_{\mathrm e}^2c^2
 &=m_{\mathrm e}^2c^2+2p'\!\cdot k,\\
 p'\!\cdot k&=0.
\end{align*}
对未来指向的时间样 $p'$ 与未来指向的类光 $k$，若 $k\neq0$ 则 $p'\!\cdot k>0$；所以唯一解是
\begin{equation*}
 {k^\mu=0.}
 \end{equation*}
也就是说，无质量-光子极限中自能分支点与电子质量壳相接，但没有有限相空间的 $e^-\to e^-\gamma$ 衰变通道。因此物理电子不获得通常 unstable-粒子意义下的有限衰变宽度。
\end{derivation}

\begin{derivation}[从\eqnrefs{eq:gordon-identity-r14}到\eqnrefs{eq:form-factor-definition-r14}]
Gordon 分解把在壳电流中的动量和结构换成 Dirac 与 Pauli 两类张量。设两条外部电子均在质量壳上，

\begin{equation*}
 p^2=p'^2=m_{\mathrm e}^2c^2,
 \qquad
 (\slashed p-m_{\mathrm e} c)u(p)=0,
 \qquad
 \bar u(p')(\slashed p'-m_{\mathrm e} c)=0,
\end{equation*}
并定义四动量转移
\begin{equation*}
 q^\mu\equiv p'^\mu-p^\mu.
\end{equation*}
左、右 Dirac equations 分别乘在 $\gamma^\mu$ 两侧：
\begin{align*}
 2m_{\mathrm e} c\,\bar u(p')\gamma^\mu u(p)
 &=\bar u(p')\left(\slashed p'\gamma^\mu+\gamma^\mu\slashed p\right)u(p).
\end{align*}
由
\begin{equation*}
 \gamma^\rho\gamma^\mu
 =\eta^{\rho\mu}+\ii\sigma^{\mu\rho},
 \qquad
 \gamma^\mu\gamma^\rho
 =\eta^{\mu\rho}-\ii\sigma^{\mu\rho},
\end{equation*}
逐项得到
\begin{align*}
 \slashed p'\gamma^\mu
 &=p'^\mu+\ii\sigma^{\mu\nu}p'_\nu,\\
 \gamma^\mu\slashed p
 &=p^\mu-\ii\sigma^{\mu\nu}p_\nu.
\end{align*}
因此
\begin{equation*}
 {
 \bar u(p')\gamma^\mu u(p)
 =\frac1{2m_{\mathrm e} c}\bar u(p')
 \left[(p'+p)^\mu+\ii\sigma^{\mu\nu}q_\nu\right]u(p).}
 \end{equation*}
这给出 Gordon 恒等式；由于 $p^\mu$ 与 $q^\mu$ 都具有动量量纲，分母必须为 $2m_{\mathrm e} c$。

对宇称偶、Lorentz 协变的电磁流，在壳旋量之间任何含 $(p'+p)^\mu$ 的结构都可用 \eqnrefs{eq:gordon-identity-r14} 改写成 $\gamma^\mu$ 与 $\sigma^{\mu\nu}q_\nu$；含 $q^\mu$ 的纵向项在守恒外部电流中由 Ward 恒等式消去。因此定义两个无量纲形状因子
\begin{equation*}
 {
 \bar u(p')\Gamma^\mu(p',p)u(p)
 =\bar u(p')\left[
 F_1(q^2)\gamma^\mu
 +\frac{\ii\sigma^{\mu\nu}q_\nu}{2m_{\mathrm e} c}F_2(q^2)
 \right]u(p).}
 \end{equation*}
$F_1$ 称 Dirac 形状因子，控制 electric-电荷/电流结构；$F_2$ 称 Pauli 形状因子，控制相对于最小耦合 Dirac 电流新增的自旋--magnetic 张量结构。二者都是无量纲函数。
\end{derivation}

\begin{derivation}[从\eqnrefs{eq:vertex-loop-start-r14}到\eqnrefs{eq:vertex-parameterized-mother-dp75}]
由\eqnrefs{eq:feynman-three-denominator-r14}，分母线性组合为
\begin{align*}
xD'_e+yD_e+zD_\gamma
&=\ell^2-2\ell\cdot(xp'+yp)
+x(p'^2-m_{\mathrm e}^2c^2)+y(p^2-m_{\mathrm e}^2c^2)+\ii0^+\\
&=\ell^2-2\ell\cdot(xp'+yp)+\ii0^+.
\end{align*}
令 $L=\ell-xp'-yp$，则
\begin{equation*}
xD'_e+yD_e+zD_\gamma
=L^2-(xp'+yp)^2+\ii0^+.
\end{equation*}
由 $q=p'-p$ 与 $p^2=p'^2=m_{\mathrm e}^2c^2$，
\begin{align*}
2p\!\cdot p'&=2m_{\mathrm e}^2c^2-q^2,\\
(xp'+yp)^2
&=(x+y)^2m_{\mathrm e}^2c^2-xyq^2\\
&=(1-z)^2m_{\mathrm e}^2c^2-xyq^2
=\Delta_V(x,y,z;q^2).
\end{align*}
把这一结果代回三分母母公式即得\eqnrefs{eq:vertex-parameterized-mother-dp75}。
\end{derivation}

\begin{derivation}[从\eqnrefs{eq:vertex-parameterized-mother-dp75}到\eqnrefs{eq:vertex-pauli-numerator-r68}]

令
\begin{equation*}
 R'^\mu\equiv p'^\mu-\ell^\mu,
 \qquad
 R^\mu\equiv p^\mu-\ell^\mu.
\end{equation*}
用正文\eqnrefs{eq:vertex-shift-delta-dp75} 的 $L^\mu=\ell^\mu-xp'^\mu-yp^\mu$ 与 $x+y+z=1$，有两组等价写法
\begin{align*}
 R'
 &=-L+zp'+yq
 =-L+zp+(1-x)q,\\
 R
 &=-L+zp-xq
 =-L+zp'-(1-y)q.
 \end{align*}
所有关于 $L$ 的奇函数在平移不变的正则化中积分为零。含两个 $L$ 的项经过张量约化
$L^\rho L^\sigma\to\eta^{\rho\sigma}L^2/d$ 后只产生 $\gamma^\mu$ 结构，因此只影响 $F_1$。Pauli 系数本身紫外有限；完成这一投影后不存在可与 $d-4$ 相乘留下有限余项的紫外极点，所以随后 Pauli 分子的 Clifford 化简可在四维进行。Pauli 结构可从 $L$ 无关分子中单独抽出。

令
\begin{equation*}
 A^\mu\equiv zp^\mu+(1-x)q^\mu,
 \qquad
 B^\mu\equiv zp'^\mu-(1-y)q^\mu,
\end{equation*}
则相关矩阵是
\begin{equation*}
 N_0^\mu
 =\gamma^\alpha(\slashed A+m_{\mathrm e} c)
 \gamma^\mu(\slashed B+m_{\mathrm e} c)\gamma_\alpha.
\end{equation*}
由 Clifford 代数得到所需的五伽马矩阵夹心恒等式。由
\begin{align*}
 \gamma^\alpha\slashed A\gamma^\mu\slashed B\gamma_\alpha
 &=2\gamma^\mu\slashed B\slashed A
 -\slashed A(\gamma^\alpha\gamma^\mu\slashed B\gamma_\alpha),
\end{align*}
而
\begin{align*}
 \gamma^\alpha\gamma^\mu\slashed B\gamma_\alpha
 &=2\slashed B\gamma^\mu
 -\gamma^\mu(\gamma^\alpha\slashed B\gamma_\alpha),\\
 &=2\slashed B\gamma^\mu+2\gamma^\mu\slashed B,\\
 &=4B^\mu\id(),
\end{align*}
所以
\begin{align*}
 \gamma^\alpha\slashed A\gamma^\mu\slashed B\gamma_\alpha
 &=2\gamma^\mu\slashed B\slashed A-4B^\mu\slashed A,\\
 &=-2\slashed B\gamma^\mu\slashed A.
 \end{align*}
最后一步使用
$\slashed B\gamma^\mu=2B^\mu-\gamma^\mu\slashed B$。其余三个收缩分别计算为
\begin{align*}
 \gamma^\alpha\gamma^\mu\slashed B\gamma_\alpha
 &=\gamma^\alpha\gamma^\mu\gamma^\beta\gamma_\alpha B_\beta
 =4\eta^{\mu\beta}B_\beta
 =4B^\mu,\\
 \gamma^\alpha\slashed A\gamma^\mu\gamma_\alpha
 &=A_\beta\gamma^\alpha\gamma^\beta\gamma^\mu\gamma_\alpha
 =4A^\mu,\\
 \gamma^\alpha\gamma^\mu\gamma_\alpha
 &=(2\eta^{\alpha\mu}-\gamma^\mu\gamma^\alpha)\gamma_\alpha
 =2\gamma^\mu-4\gamma^\mu
 =-2\gamma^\mu.
\end{align*}
于是
\begin{equation*}
 N_0^\mu
 =-2\slashed B\gamma^\mu\slashed A
 +4m_{\mathrm e} c(A+B)^\mu
 -2m_{\mathrm e}^2c^2\gamma^\mu.
 \end{equation*}

把它夹在 $\bar u(p')$ 与 $u(p)$ 之间，并把可用的在壳 Dirac 方程移到最外侧，使用
\begin{equation*}
 \slashed B=z\slashed p'-(1-y)\slashed q,
 \qquad
 \slashed A=z\slashed p+(1-x)\slashed q.
\end{equation*}
因此
\begin{align*}
 \bar u\slashed B\gamma^\mu\slashed A u
 =&\,\bar u\Bigl[
 z^2m_{\mathrm e}^2c^2\gamma^\mu
 +zm_{\mathrm e} c(1-x)\gamma^\mu\slashed q\\
 &\qquad-zm_{\mathrm e} c(1-y)\slashed q\gamma^\mu
 -(1-x)(1-y)\slashed q\gamma^\mu\slashed q
 \Bigr]u.
\end{align*}
利用
\begin{align*}
 \gamma^\mu\slashed q
 &=q^\mu-\ii\sigma^{\mu\nu}q_\nu,\\
 \slashed q\gamma^\mu
 &=q^\mu+\ii\sigma^{\mu\nu}q_\nu,
\end{align*}
线性 $q$ 的两项给
\begin{equation*}
 zm_{\mathrm e} c\left[(y-x)q^\mu
 -\ii(2-x-y)\sigma^{\mu\nu}q_\nu\right]
 =zm_{\mathrm e} c\left[(y-x)q^\mu
 -\ii(1+z)\sigma^{\mu\nu}q_\nu\right].
\end{equation*}
而
\begin{equation*}
 \slashed q\gamma^\mu\slashed q
 =2q^\mu\slashed q-q^2\gamma^\mu
\end{equation*}
在在壳旋量夹心中满足
\begin{align*}
 \bar u(p')\slashed q\,u(p)
 &=\bar u(p')(\slashed p'-\slashed p)u(p),\\
 &=m_{\mathrm e} c\,\bar uu-m_{\mathrm e} c\,\bar uu=0,
\end{align*}
故该项只进入 $\gamma^\mu$，不贡献 Pauli 结构。

再处理 $N_0^\mu$ 的第二项：
\begin{equation*}
 A+B=z(p'+p)+(y-x)q.
\end{equation*}
其中 $(y-x)q^\mu$ 在对称 simplex 积分 $x\leftrightarrow y$ 下为奇函数，积分为零。对 $(p'+p)^\mu$ 使用 Gordon 恒等\eqnrefs{eq:gordon-identity-r14}，
\begin{equation*}
 (p'+p)^\mu
 \doteq 2m_{\mathrm e} c\gamma^\mu-\ii\sigma^{\mu\nu}q_\nu,
\end{equation*}
其中 $\doteq$ 表示两侧夹在同一对在壳旋量之间时相等。因此，第一项
$-2\slashed B\gamma^\mu\slashed A$ 的 Pauli 部分是
\begin{equation*}
 +2\ii zm_{\mathrm e} c(1+z)\sigma^{\mu\nu}q_\nu,
\end{equation*}
第二项 $4m_{\mathrm e} c(A+B)^\mu$ 的 Pauli 部分是
\begin{equation*}
 -4\ii zm_{\mathrm e} c\sigma^{\mu\nu}q_\nu.
\end{equation*}
相加得到
\begin{equation*}
 N_0^\mu\Big|_{\rm Pauli}
 =-2\ii z(1-z)m_{\mathrm e} c\,\sigma^{\mu\nu}q_\nu.
\end{equation*}
按 \eqnrefs{eq:form-factor-definition-r14} 的标准基底写成
\begin{equation*}
 N_0^\mu\Big|_{\rm Pauli}
 =-N_2\frac{\ii\sigma^{\mu\nu}q_\nu}{2m_{\mathrm e} c},
\end{equation*}
便得到
\begin{equation*}
 {
 N_2(x,y,z)=4z(1-z)m_{\mathrm e}^2c^2.}
 \end{equation*}
该结果不依赖圈动量，因此 $F_2$ 的动量积分比 $F_1$ 更收敛。
\end{derivation}

\begin{derivation}[从\eqnrefs{eq:vertex-parameterized-mother-dp75}到\eqnrefs{eq:schwinger-F2-r14}]

由参数化后的顶角积分和刚得到的 $N_2=4z(1-z)m_{\mathrm e}^2c^2$，一圈 Pauli 形状因子是
\begin{equation*}
 F_2^{\rm 1L}(q^2)
 =2\ii g_{\rm em}^2
 \int_0^1\dd x\dd y\dd z\,
 \delta(x+y+z-1)
 \int\frac{\dd^4 L}{(2\pi)^4}
 \frac{N_2}{(L^2-\Delta_V+\ii0^+)^3}.
 \end{equation*}
Pauli 投影对应的积分已经紫外有限；完成正则化代数后可以取 $d=4$。

对其中的标量圈积分直接作 Wick 转动 $L^0=\ii L_E^0$。由于 $\dd^4 L=\ii\dd^4 L_E$ 且 $L^2-\Delta=-(L_E^2+\Delta)$，有
\begin{equation*}
 \int\frac{\dd^4 L}{(2\pi)^4}\frac1{(L^2-\Delta+\ii0^+)^3}
 =-\ii\int\frac{\dd^4 L_E}{(2\pi)^4}\frac1{(L_E^2+\Delta)^3}.
\end{equation*}
按正文的 $\Omega_d\equiv\int\dd\Omega_{d-1}$ 记号，四维 Euclid 空间的单位三球面面积为 $\Omega_4=2\pi^2$；写 $R=|L_E|$，
\begin{align*}
 \int\frac{\dd^4 L}{(2\pi)^4}\frac1{(L^2-\Delta+\ii0^+)^3}
 &=-\ii\frac{2\pi^2}{(2\pi)^4}
 \int_0^\infty\dd R\,
 \frac{R^3}{(R^2+\Delta)^3},\\
 &=-\frac{\ii}{8\pi^2}\frac12
 \int_0^\infty\dd u\,
 \frac{u}{(u+\Delta)^3},
 \qquad u\equiv R^2.
\end{align*}
再令 $u=\Delta v$：
\begin{align*}
 \int_0^\infty\dd u\frac{u}{(u+\Delta)^3}
 &=\frac1\Delta
 \int_0^\infty\dd v\frac{v}{(1+v)^3},\\
 &=\frac1\Delta
 \left[\frac12\right],
\end{align*}
其中最后一项可由 $w=1/(1+v)$ 直接积分得到。因此
\begin{equation*}
 {
 \int\frac{\dd^4 L}{(2\pi)^4}\frac1{(L^2-\Delta+\ii0^+)^3}
 =-\frac{\ii}{32\pi^2}\frac1\Delta.}
 \end{equation*}

现在取 $q^2=0$。由参数化分母
\begin{equation*}
 \Delta_V=(1-z)^2m_{\mathrm e}^2c^2,
\end{equation*}
所以
\begin{equation*}
 \frac{N_2}{\Delta_V}
 =\frac{4z(1-z)m_{\mathrm e}^2c^2}{(1-z)^2m_{\mathrm e}^2c^2}
 =\frac{4z}{1-z}.
\end{equation*}
把上面求得的标量圈积分代入 Pauli 形状因子的母积分：
\begin{align*}
 F_2^{\rm 1L}(0)
 &=\frac{g_{\rm em}^2}{16\pi^2}
 \int_0^1\dd x\dd y\dd z\,
 \delta(x+y+z-1)
 \frac{4z}{1-z}.
\end{align*}
固定 $z$ 后，单纯形上 $x+y=1-z$，故
\begin{align*}
 \int_0^1\dd x\dd y\,
 \delta(x+y+z-1)
 &=\int_0^{1-z}\dd x=1-z.
\end{align*}
该几何因子恰好消去被积函数中的 $1-z$：
\begin{align*}
 F_2^{\rm 1L}(0)
 &=\frac{g_{\rm em}^2}{16\pi^2}
 \int_0^1\dd z\,4z,\\
 &=\frac{g_{\rm em}^2}{8\pi^2},\\
 &={\frac{\alpha}{2\pi}},
 \end{align*}
最后一步用了 $g_{\rm em}^2=4\pi\alpha$。整个计算中 $m_{\mathrm e}^2c^2$ 在分子与分母中显式相消，因此 $F_2(0)$ 自然保持无量纲。
\end{derivation}

\begin{derivation}[从\eqnrefs{eq:schwinger-F2-r14}到\eqnrefs{eq:ae-form-factor-r68}]
最小耦合 Dirac Hamilton 量的低能展开给出自旋磁矩项

\begin{equation*}
 H_{B}^{\rm Dirac}
 =-\frac{q_{\mathrm e}\hbar}{2m_{\mathrm e}}\boldsymbol\sigma_{\rm P}\cdot\mathbf B
 =-\frac{q_{\mathrm e}}{m_{\mathrm e}}\mathbf S_{1/2}\cdot\mathbf B,
\end{equation*}
其中 $\mathbf S_{1/2}=\hbar\boldsymbol\sigma_{\rm P}/2$。正文已经用
$H_B=-(g_{\rm mag}q_{\mathrm e}/2m_{\mathrm e})\mathbf S_{1/2}\cdot\mathbf B$ 固定了旋磁因子的物理含义；这里仅比较顶角中两个磁张量的系数。

由 Gordon 恒等式，$F_1\gamma^\mu$ 自身已经含有
$F_1\,\ii\sigma^{\mu\nu}q_\nu/(2m_{\mathrm e} c)$ 磁张量；Pauli 形状因子又增加
$F_2\,\ii\sigma^{\mu\nu}q_\nu/(2m_{\mathrm e} c)$。因此静态极限 $q^2\to0$ 的磁耦合相对于树级 Dirac 结果乘上
$F_1(0)+F_2(0)$：
\begin{equation*}
 {
 g_{\rm mag}=2\,[F_1(0)+F_2(0)].}
 \end{equation*}
物理电荷归一化与 Ward--Takahashi 恒等式共同固定
\begin{equation*}
 F_1(0)=1.
\end{equation*}
正文\eqnrefs{eq:ae-form-factor-r68} 已把相对 Dirac 值的偏移记为 $a_\ee^{\rm anom}=(g_{\rm mag}-2)/2$，因此
\begin{equation*}
 a_\ee^{\rm anom}=F_2(0).
\end{equation*}
结合 \eqnrefs{eq:schwinger-F2-r14} 得到 Schwinger 领头修正
\begin{equation*}
 {
 a_\ee^{\rm anom}
 =\frac{\alpha}{2\pi}+\order(\alpha^2),
 \qquad
 g_{\rm mag}
 =2+\frac{\alpha}{\pi}+\order(\alpha^2).}
 \end{equation*}
\end{derivation}

\begin{derivation}[从\eqnrefs{eq:vertex-parameterized-mother-dp75}到\eqnrefs{eq:F2-spacelike-ratio-r14}]

同一计算还可以保留 $q^2\neq0$。正文已经取空间样转移动量 $Q^2=-q^2>0$；在国际单位制中 $Q$ 的量纲是动量。代入上面已经直接算出的标量圈积分，得到
\begin{equation*}
 F_2^{\rm 1L}(q^2)
 =\frac{g_{\rm em}^2}{16\pi^2}
 \int_0^1\dd x\dd y\dd z\,
 \delta(x+y+z-1)
 \frac{4z(1-z)m_{\mathrm e}^2c^2}
 {(1-z)^2m_{\mathrm e}^2c^2-xyq^2}.
 \end{equation*}
令
\begin{equation*}
 w\equiv1-z=x+y,
 \qquad
 u_F\equiv\frac{x}{x+y},
 \qquad
 x=w u_F,
 \qquad
 y=w(1-u_F).
\end{equation*}
这里 $w,u_F\in[0,1]$ 都是无量纲变量。雅可比因子是
\begin{equation*}
 \left|\det\frac{\partial(x,y)}{\partial(w,u_F)}\right|=w.
\end{equation*}
于是分母与分子分别变成
\begin{align*}
 (1-z)^2m_{\mathrm e}^2c^2-xyq^2
 &=w^2[m_{\mathrm e}^2c^2-u_F(1-u_F)q^2],\\
 4z(1-z)m_{\mathrm e}^2c^2
 &=4(1-w)w\,m_{\mathrm e}^2c^2.
\end{align*}
连同测度的额外 $w$，所有 $w^2$ 相消：
\begin{align*}
 F_2^{\rm 1L}(q^2)
 &=\frac{g_{\rm em}^2}{16\pi^2}
 \int_0^1\dd u_F
 \frac{m_{\mathrm e}^2c^2}{m_{\mathrm e}^2c^2-u_F(1-u_F)q^2}
 \int_0^1\dd w\,4(1-w),\\
 &={
 \frac{\alpha}{2\pi}
 \int_0^1\dd u_F\,
 \frac{m_{\mathrm e}^2c^2}
 {m_{\mathrm e}^2c^2-u_F(1-u_F)q^2}}.
 \end{align*}
在空间样区域 $q^2=-Q^2$，被积函数没有极点：
\begin{equation*}
 \frac{F_2^{\rm 1L}(-Q^2)}{F_2^{\rm 1L}(0)}
 =\int_0^1\dd u_F\,
 \frac1{1+[Q/(m_{\mathrm e} c)]^2u_F(1-u_F)}.
 \end{equation*}
\end{derivation}

\begin{derivation}[从\eqnrefs{eq:ImPi-r12}到\eqnrefs{eq:Uehling-momentum-correction-r12}]
取点状原子核电荷为 $+Ze$、电子电荷为 $q_{\mathrm e}=-e$。静态 Coulomb 势能为

\begin{equation*}
 V_C(r)=-\frac{Ze^2}{4\pi\varepsilon_0r}
 =-\frac{Z\alpha\hbar c}{r}.
 \end{equation*}
空间 Fourier 变换写为
\begin{equation*}
 V(\mathbf r)
 =\int\frac{\dd^3 q}{(2\pi\hbar)^3}
 \ee^{+\ii\mathbf q\cdot\mathbf r/\hbar}
 \widetilde V(\mathbf q).
 \end{equation*}
由
\begin{equation*}
 \int\frac{\dd^3 q}{(2\pi\hbar)^3}
 \ee^{+\ii\mathbf q\cdot\mathbf r/\hbar}
 \frac{\hbar^2}{\mathbf q^2}
 =\frac1{4\pi r},
\end{equation*}
可知
\begin{equation*}
 \widetilde V_C(\mathbf q)
 =-\frac{Ze^2\hbar^2}{\varepsilon_0\mathbf q^2}.
 \end{equation*}

静态交换满足 $q^0=0$，因而 $q^2=-Q^2$，其中 $Q\equiv|\mathbf q|$。按真空极化的符号约定，定义空间样屏蔽函数
\begin{equation*}
 \widehat\Pi(Q^2)
 \equiv-\Pi_R(-Q^2)
 =\frac{2\alpha}{\pi}
 \int_0^1\dd x\,x(1-x)
 \ln\left[1+\frac{Q^2}{m_{\mathrm e}^2c^2}x(1-x)\right].
 \end{equation*}
Dyson-resummed 横向光子分母是 $Q^2[1+\Pi_R(-Q^2)]$。保留一圈的一阶修正，
\begin{equation*}
 \frac1{Q^2[1+\Pi_R(-Q^2)]}
 =\frac1{Q^2}[1+\widehat\Pi(Q^2)]+
 \order(\alpha^2).
\end{equation*}
因此
\begin{equation*}
 \delta\widetilde V_U(\mathbf q)
 =\widetilde V_C(\mathbf q)\widehat\Pi(Q^2).
 \end{equation*}

三维 Fourier 变换可借助一次减除的色散表示。$\Pi_R(z)/z$ 在复 $z$ 平面除 $z\ge4m_{\mathrm e}^2c^2$ 的支割外解析，并且 $\Pi_R(0)=0$。对包围支割的 Cauchy 围道，
\begin{equation*}
 \frac{\Pi_R(z)}{z}
 =\frac1{2\pi\ii}
 \int_{4m_{\mathrm e}^2c^2}^{\infty}\dd s\,
 \frac{\Pi_R(s+\ii0)-\Pi_R(s-\ii0)}{s(s-z)}.
\end{equation*}
由于不连续度为 $2\ii\operatorname{Im}\Pi_R(s)$，取 $z=-Q^2$ 得
\begin{equation*}
 \widehat\Pi(Q^2)
 =\frac{Q^2}{\pi}
 \int_{4m_{\mathrm e}^2c^2}^{\infty}\dd s\,
 \frac{\operatorname{Im}\Pi_R(s)}{s(s+Q^2)}.
 \end{equation*}
再代入 \eqnrefs{eq:ImPi-r12}：
\begin{equation*}
 \frac{\widehat\Pi(Q^2)}{Q^2}
 =\frac{\alpha}{3\pi}
 \int_{4m_{\mathrm e}^2c^2}^{\infty}\frac{\dd s}{s(s+Q^2)}
 \sqrt{1-\frac{4m_{\mathrm e}^2c^2}{s}}
 \left(1+\frac{2m_{\mathrm e}^2c^2}{s}\right).
 \end{equation*}

相应的 Fourier 变换为
\begin{equation*}
 I_s(r)
 \equiv
 \int\frac{\dd^3 q}{(2\pi\hbar)^3}
 \ee^{+\ii\mathbf q\cdot\mathbf r/\hbar}
 \frac{\hbar^2}{\mathbf q^2+s}.
\end{equation*}
取球坐标并把极轴沿 $\mathbf r$，角积分为
\begin{equation*}
 \int\dd\Omega_q\,\ee^{\ii Qr\cos\theta/\hbar}
 =4\pi\frac{\sin(Qr/\hbar)}{Qr/\hbar}.
\end{equation*}
所以
\begin{equation*}
 I_s(r)
 =\frac1{2\pi^2r}
 \int_0^\infty\dd Q\,
 \frac{Q\sin(Qr/\hbar)}{Q^2+s}.
\end{equation*}
将 $\sin$ 写成指数并在上半平面闭合积分轮廓，唯一极点在 $Q=+\ii\sqrt s$；留数定理给
\begin{equation*}
 \int_0^\infty\dd Q\,
 \frac{Q\sin(Qr/\hbar)}{Q^2+s}
 =\frac\pi2\ee^{-\sqrt s\,r/\hbar},
\end{equation*}
因此
\begin{equation*}
 {
 I_s(r)=\frac{\ee^{-\sqrt s\,r/\hbar}}{4\pi r}.}
 \end{equation*}
把 \eqnrefs{eq:Pi-dispersion-explicit-r12} 插入 \eqnrefs{eq:Uehling-momentum-correction-r12}，交换 $s$ 与 $\mathbf q$ 积分，并使用上面得到的三维 Fourier 积分：
\begin{align*}
 \delta V_U(r)
 &=V_C(r)\frac{\alpha}{3\pi}
 \int_{4m_{\mathrm e}^2c^2}^{\infty}\frac{\dd s}{s}
 \ee^{-\sqrt s\,r/\hbar}
 \sqrt{1-\frac{4m_{\mathrm e}^2c^2}{s}}
 \left(1+\frac{2m_{\mathrm e}^2c^2}{s}\right).
\end{align*}
令
\begin{equation*}
 \sqrt s=2m_{\mathrm e} c\,t,
 \qquad t\in[1,\infty),
 \qquad \frac{\dd s}{s}=\frac{2\dd t}{t},
\end{equation*}
并定义约化 Compton 波长 $\lambda_C\equiv\hbar/(m_{\mathrm e} c)$，得到经典 Uehling 修正：
\begin{equation*}
 {
 \delta V_U(r)
 =V_C(r)\frac{2\alpha}{3\pi}
 \int_1^\infty\dd t\,
 \ee^{-2rt/\lambda_C}
 \left(1+\frac1{2t^2}\right)
 \frac{\sqrt{t^2-1}}{t^2}.}
 \end{equation*}
\end{derivation}

\begin{derivation}[从\eqnrefs{eq:Volkov-Dirac-r12}到\eqnrefs{eq:Volkov-solution-r12}]
强经典激光背景不应在一开始按光子数截断。令经典四势只依赖平面波相位

\begin{equation*}
 \phi\equiv\kappa\cdot x,
 \qquad
 \kappa^2=0,
 \qquad
 \kappa\cdot A(\phi)=0,
 \end{equation*}
其中 $\kappa^\mu=(\omega_L/c,\mathbf k_L)$ 是四波矢，量纲为 $\mathrm{m^{-1}}$；$p^\mu$ 仍是电子物理四动量。Dirac 方程写成
\begin{equation*}
 [\ii\hbar\slashed\partial-q_{\mathrm e}\slashed A(\phi)-m_{\mathrm e} c]\psi(x)=0.
 \end{equation*}
原式整体除以 $c$ 后，三项都具有动量量纲。

取自由在壳自旋子 $u_s(p)$，满足
\begin{equation*}
 (\slashed p-m_{\mathrm e} c)u_s(p)=0.
\end{equation*}
定义
\begin{align*}
 R_p(\phi)
 &\equiv
 1+\frac{q_{\mathrm e}}{2p\cdot\kappa}
 \slashed\kappa\slashed A(\phi),\\
 F_p(\phi)
 &\equiv
 \frac{q_{\mathrm e} p\cdot A(\phi)}{p\cdot\kappa}
 -\frac{q_{\mathrm e}^2A^2(\phi)}{2p\cdot\kappa},\\
 S_p(x)
 &\equiv p\cdot x+\int^{\phi}\dd\varphi\,F_p(\varphi).
 \end{align*}
量纲检查很直接：$p\cdot\kappa$ 是动量/长度，$q_{\mathrm e} A$ 是动量，因此 $R_p$ 无量纲；$F_p$ 具有作用量量纲，因为相位变量 $\phi$ 无量纲，所以 $S_p/\hbar$ 无量纲。

候选解为
\begin{equation*}
 {
 \psi_{p,s}^{\rm V}(x)
 =R_p(\phi)u_s(p)
 \exp\left[-\frac{\ii}{\hbar}S_p(x)\right].}
 \end{equation*}
将该形式直接代回 Dirac 方程。首先
\begin{equation*}
 \partial_\mu S_p=p_\mu+\kappa_\mu F_p,
 \qquad
 \partial_\mu R_p=\kappa_\mu R_p'(\phi).
\end{equation*}
由于 $\slashed\kappa^2=\kappa^2=0$，
\begin{equation*}
 \slashed\partial R_p
 =\slashed\kappa R_p'
 =\frac{q_{\mathrm e}}{2p\cdot\kappa}
 \slashed\kappa\slashed\kappa\slashed A'(\phi)
 =0.
 \end{equation*}
所以\eqnrefs{eq:Volkov-Dirac-r12} 作用在候选解上只剩
\begin{equation*}
 [\slashed p+\slashed\kappa F_p-q_{\mathrm e}\slashed A-m_{\mathrm e} c]R_pu_s.
 \end{equation*}
需要证明它为零。先计算 $\slashed pR_pu_s$。利用
\begin{equation*}
 \slashed p\slashed\kappa=2p\cdot\kappa-\slashed\kappa\slashed p,
 \qquad
 \slashed p\slashed A=2p\cdot A-\slashed A\slashed p,
\end{equation*}
以及 $\slashed p u_s=m_{\mathrm e} cu_s$：
\begin{align*}
 \slashed pR_pu_s
 &=m_{\mathrm e} cu_s+\frac{q_{\mathrm e}}{2p\cdot\kappa}
 \slashed p\slashed\kappa\slashed A u_s,\\
 &=m_{\mathrm e} cu_s+q_{\mathrm e}\slashed A u_s
 -\frac{q_{\mathrm e}}{2p\cdot\kappa}
 \slashed\kappa\slashed p\slashed A u_s,\\
 &=m_{\mathrm e} cR_pu_s+q_{\mathrm e}\slashed A u_s
 -q_{\mathrm e}\frac{p\cdot A}{p\cdot\kappa}\slashed\kappa u_s.
 \end{align*}
另一方面，$\kappa\cdot A=0$ 给
\begin{equation*}
 \slashed A\slashed\kappa\slashed A
 =-A^2\slashed\kappa,
\end{equation*}
于是
\begin{equation*}
 q_{\mathrm e}\slashed A R_pu_s
 =q_{\mathrm e}\slashed A u_s
 -\frac{q_{\mathrm e}^2A^2}{2p\cdot\kappa}\slashed\kappa u_s.
 \end{equation*}
同时 $\slashed\kappa R_p=\slashed\kappa$，故
\begin{equation*}
 \slashed\kappa F_p R_pu_s
 =\left[
 q_{\mathrm e}\frac{p\cdot A}{p\cdot\kappa}
 -\frac{q_{\mathrm e}^2A^2}{2p\cdot\kappa}
 \right]\slashed\kappa u_s.
 \end{equation*}
$\slashed pR_pu_s$、$q_{\mathrm e}\slashed A R_pu_s$ 与 $\slashed\kappa F_pR_pu_s$ 三个结果相加后，$q_{\mathrm e}\slashed A$、$p\cdot A$ 与 $A^2$ 三类项逐项相消，最后只剩
\begin{equation*}
 m_{\mathrm e} cR_pu_s-m_{\mathrm e} cR_pu_s=0.
\end{equation*}
因此 \eqnrefs{eq:Volkov-solution-r12} 确实解含任意强度经典平面波的 Dirac 方程。

对周期平面波且选规范使 $\langle A\rangle_\phi=0$，长时间平均相位中的线性项消失。定义 Volkov 准动量
\begin{equation*}
 {
 \bar p^\mu
 =p^\mu
 -\frac{q_{\mathrm e}^2\langle A^2\rangle_\phi}{2p\cdot\kappa}\kappa^\mu.}
 \end{equation*}
因为 $\kappa^2=0$，
\begin{align*}
 \bar p^2
 &=p^2-q_{\mathrm e}^2\langle A^2\rangle_\phi,\\
 &=m_{\mathrm e}^2c^2-q_{\mathrm e}^2\langle A^2\rangle_\phi.
\end{align*}
横向平面波有 $A^2<0$，所以准动量对应的有效质量增大。若定义峰值归一化矢势
\begin{equation*}
 a_0\equiv\frac{|q_{\mathrm e}|A_0}{m_{\mathrm e} c},
\end{equation*}
则圆偏振的 $\langle A^2\rangle=-A_0^2$ 给
\begin{equation*}
 {m_*^{\rm circ}=m_{\mathrm e}\sqrt{1+a_0^2},}
\end{equation*}
而线偏振可取 $A^\mu(\phi)=\mathcal A^\mu\cos\phi$，并定义正的峰值幅度 $A_0\equiv\sqrt{-\mathcal A^2}$；于是 $\langle A^2\rangle=\mathcal A^2/2=-A_0^2/2$，故
\begin{equation*}
 {m_*^{\rm lin}=m_{\mathrm e}\sqrt{1+\frac{a_0^2}{2}}.}
 \end{equation*}
这些是完整平面波 Dirac 动力学的结果；只有在随后再取 $a_0\ll1$ 时，才可以把它展开成普通的弱场多光子微扰绘景。
\end{derivation}
